\documentclass[12pt,a4paper,oneside]{article}
\usepackage[brazil]{babel}
\usepackage[utf8]{inputenc}
\usepackage{olymp}

\setcounter{problem}{3}

\begin{document}

\begin{problem}{Risada automática}{risada}{0}
 
Cansado de ter que comentar piadas e vídeos enviados em grupos de família e escola, Joaquim passou a responder apenas \texttt{kkkkk}, com diferente números de {k}'s dependendo de quão engraçado (ou sem graça) for o que tiverem postado. Mas ele também já se cansou de digitar tantos {k}'s e pediu sua ajuda para automatizar suas risadas. 


%\Notes
%
%Observações, se tiver.

\InputFile

A entrada contém um único número natural $N$. Restrições: $0 < N \leq 30$.


\OutputFile

Seu programa deve gerar apenas uma linha de saída, contendo uma sequência de $N$ \texttt{k}'s.

% \Notes
% Preste atenção aos tipos das variáveis, pois $20! \approx 2^{61}$.

\Example

\begin{example}
\exmp{
5
}{
kkkkk
}%
\end{example}

\begin{example}
\exmp{
3
}{
kkk
}%
\end{example}

\begin{example}
\exmp{
10
}{
kkkkkkkkkk
}%
\end{example}



%Comentários sobre o exemplo, se necessário.

\end{problem}

\end{document}
